\section{Conductor products and coefficient-prime scope}\label{sec:conductors}
For a realization $V_\ell$ and a prime $p\ne\ell$, use
\[
 P_p(V_\ell,T)=\det(1-\Frob_pT\mid V_\ell^{I_p}),\qquad
 a_p(V_\ell)=\dim V_\ell-\dim V_\ell^{I_p}+\Sw_p(V_\ell).
\]
The local proofs are carried out directly with arbitrary coefficient prime
different from the residue characteristic. They give
\[
\begin{array}{c|ccc|ccc}
 &d_2&\Sw_2&a_2&d_3&\Sw_3&a_3\\\hline
 W_{4,\ell}^+&0&0&4&1&0&3\\
 W_{4,\ell}^-&0&0&4&0&0&4
\end{array}
\]
with the $p$-column interpreted only for $\ell\ne p$. At $3$ the invariant
line of the plus descent has geometric Frobenius eigenvalue $3$.
Thus $P_2^+=P_2^-=P_3^-=1$ and $P_3^+=1-3T$.

For $p\ge5$ and $\ell\ne p$, good-prime unramifiedness follows by the
relative tame-specialization argument of Section~\ref{sec:geometry}.
To justify its coefficient change explicitly, use the integral sheaf on
the full cubic cover before applying rational projectors. At boundary points of both generic and special geometric fibres,
the Tate-module description gives prime-to-$p$ unipotent local
monodromy; at infinity the additional quadratic twist is also tame.
The cubic Kummer cover is tame since $p\ne3$. Thus every finite
$\ell$-power coefficient level is relatively tame. The same specialization,
duality and parabolic-image comparison then applies, and rational
projectors pass to its summands. This uses the local Tate descriptions,
not a claim that the order of $\GL_2(\Z/\ell^n)$ is prime to $p$.
Hence, for every $\ell\ne2,3$,
\[
 N^{(\ell)}(W_{4,\ell}^+)=\prod_{p\ne\ell}p^{a_p}=2^4\,3^3=432,
 \qquad
 N^{(\ell)}(W_{4,\ell}^-)=2^4\,3^4=1296.
\]
Equivalently these are the conductor products obtained by choosing an
auxiliary coefficient prime away from each residue characteristic.
Independence of the displayed bad-prime factors and exponents follows
from the explicit calculations, without invoking a general bad-prime
compatibility theorem. No assertion about the usual inertia invariants
of the $p$-adic realization at $p=\ell$ is being substituted for an
away-from-$\ell$ Euler factor. Nor does this calculation identify the
entire Weil--Deligne representation in all coefficient characteristics.

The numerical conductors formerly proposed for the two recipes therefore
agree with the conductor products of the specified descents. This does
not establish the functional-equation sign, analytic continuation, or
the accuracy of the historical central-value computation.
